\chapter{شروط المجتهد والمفتي وأنواع القياس}

\section{شروط الاجتهاد والفتوى}
الحمد لله والصلاة والسلام على رسول الله، أما بعد؛ فقد عقد الإمام الشافعي رحمه الله تعالى باباً نفيساً في "شروط المجتهد" ومن يحق له أن يتصدر للفتوى والقياس. وقد قال شيخ شيوخنا العلامة أحمد شاكر رحمه الله عن هذه الفِقرات: «هذه الدرر الغالية والحكم البالغة هي أحسن ما قرأت في شروط الاجتهاد».

قال الشافعي رحمه الله: «ولا ينبغي للمفتي أن يفتي أحداً إلا أن يكون عالماً بكتاب الله، وعلم ناسخه ومنسوخه، وخاصه وعامه، وأدبه. وعالماً بسنن رسول الله ﷺ، وأقاويل أهل العلم قديماً وحديثاً (مواطن الإجماع والاختلاف). وعالماً بلسان العرب، عاقلاً يميز بين المشتبه ويعقل القياس».

فإن عُدم واحدة من هذه الخصال، لم يحل له أن يقول بقياس أو يفتي. ومثّل لذلك بمن يفتقد أداة من أدوات الفهم؛ كمن كان عالماً بالأصول ولكنه لا يعقل القياس، أو كان عاقلاً للقياس ومضيعاً لعلم الأصول؛ فكلاهما لا يجوز له أن يقيس. 
كما لا يجوز أن يوصف الطريق لأعمى فيقال له: «اجعل كذا عن يمينك وكذا عن يسارك»، وهو لا يبصر. وكما لا يجوز لجاهل بأسعار السوق أن يُقيّم سلعة، فكذلك لا يقاس في حلال الله وحرامه إلا بآلة متكاملة من العلم.

\subsection{الإنصاف وأدب المجتهد في الخلاف}
من أعظم شروط المجتهد التي نص عليها الشافعي: بلوغ غاية الجهد في البحث عن الحق، والإنصاف من نفسه.
فيجب على طالب العلم ألا يمتنع من الاستماع لمن خالفه؛ لأنه قد يتنبه بالاستماع إلى ترك الغفلة، ويزداد تثبيتاً فيما اعتقده من الصواب. فالمخالف إما أن يزيدك يقيناً على ما أنت عليه إذا ظهر لك ضعف دليله، وإما أن يردك إلى الصواب إذا كنت مخطئاً. والواجب أن تنصف من نفسك فتتبع الحق حيث ظهر. ولا يكون العالم أعمى مقلداً يقلد بغير حجة.

\section{القياس ومراتبه}
بعد أن بيّن الشافعي شروط القياس والمقيس، شرع في بيان أقسام القياس ومراتبه قوةً ووضوحاً:

\subsection{أقوى القياس (قياس الأولى)}
يقرر الشافعي أن أقوى أنواع القياس هو: أن يحرّم الله في كتابه أو يحرّم رسوله ﷺ القليل من الشيء، فيُعلم يقيناً أن الكثير منه أشد تحريماً، بفضل الكثرة على القلة. 
وكذلك إذا حُمد العبد على يسير من الطاعة، كان ما هو أكثر منها أولى أن يُحمد عليه. وإذا أبيح الكثير من الشيء، كان الأقل منه أولى أن يكون مباحاً.

\textbf{أمثلة وتطبيقات:}
\begin{itemize}
    \item قال تعالى: ﴿فَمَن يَعْمَلْ مِثْقَالَ ذَرَّةٍ خَيْرًا يَرَهُ * وَمَن يَعْمَلْ مِثْقَالَ ذَرَّةٍ شَرًّا يَرَهُ﴾. فما هو أكثر من مثقال ذرة من الخير أعظم حمداً وأجراً، وما هو أكثر من ذرة من الشر أعظم مأثماً.
    \item قال ﷺ: «إن الله حرم من المؤمن دمه وماله وأن يُظن به إلا خيراً». فإذا حرم مجرد سوء الظن المخفي، فالتصريح بالقول الباطل فيه والإيذاء الفعلي أشد تحريماً من باب أولى.
    \item أباح الله لنا دماء أهل الكفر المحاربين غير المعاهدين وأموالهم بأسرها، فما كان من نيل شيء من أبدانهم دون الدماء أو أخذ بعض أموالهم فهو أولى بالإباحة.
\end{itemize}

\section{أمثلة فقهية دقيقة على القياس (العلة والشبه)}

\subsection{القياس في وجوب نفقة الأب العاجز}
قال تعالى: ﴿وَالْوَالِدَاتُ يُرْضِعْنَ أَوْلَادَهُنَّ حَوْلَيْنِ كَامِلَيْنِ ۖ لِمَنْ أَرَادَ أَن يُتِمَّ الرَّضَاعَةَ ۚ وَعَلَى الْمَوْلُودِ لَهُ رِزْقُهُنَّ وَكِسْوَتُهُنَّ بِالْمَعْرُوفِ﴾. 
دلت الآية والسنة (كحديث هند بنت عتبة: «خذي من ماله ما يكفيك وولدك بالمعروف») على أن نفقة الولد الصغير واجبة على الوالد، والعلة في ذلك أن الولد بضعة من أبيه، وأنه عاجز في حال صغره عن الكسب وإغناء نفسه. 

فقاس الشافعي على ذلك: أنه إذا بلغ الأب من الكِبَر والعجز مبلغا لا يغني فيه نفسه بكسب ولا مال، وولده غني قادر؛ وجبت النفقة لِلأب على ولده. لأن الولد من الوالد، فكما أُجبر الأب على نفقة ابنه الصغير لعجزه، يُجبر الابن على نفقة أبيه لعجزه؛ فلا يُضيّع أحدهما ما هو بضعة منه.

\subsection{القياس في قاعدة الخراج بالضمان}
قضى النبي ﷺ في رجل اشترى عبداً فاستغله (أي جعله يعمل ويُدِرُّ عليه مالاً)، ثم اكتشف فيه المشتري عيباً قديماً فرده على بائعه الأول، فطلب البائع غلة العبد (أجرته في تلك الأيام). فقضى النبي ﷺ بأن العبد يُرد بالعيب، وأن للمشتري أن يحبس الغلة لنفسه ولا يردها؛ لأنه كان «ضامناً» للعبد، فلو مات في تلك الفترة لمات من مال المشتري، فقوبل الضمان باستحقاق الخراج (الخراج بالضمان).

قاس الشافعي على هذه السنة الجليلة أشياء أخرى لم تُنص بعينها:
فمن اشترى نخلة فأثمرت، أو جارية فولدت، أو ماشية فحلبت، ثم ظهر عيب قديم فردها المشتري؛ فإن الثمر والولد واللبن من حق المشتري، لأن كل هذا حدث في فترة ملكه وضمانه، فلحق بالأصل المُقاس عليه وهو غلة العبد.

\subsection{القياس في الربويات}
نهى النبي ﷺ عن التفاضل في الذهب بالذهب، والفضة بالفضة، والتمر والبر والشعير. 
والعلة في المأكولات الثلاثة أنها مأكولة تُكال وتُدخر للاقتيات. فقاس الشافعي والجمهور عليها: العسل، والسمن، والزيت، والسكر، وغيرها مما يقاس بالوزن أو الكيل ويكون مأكولاً مدخراً. فهذه لا تُباع بجنسها إلا متماثلة ويداً بيد قياساً على المنصوص عليه. ومنع قياس المأكولات على الذهب والفضة (الدنانير والدراهم)؛ لاختلاف العلة والجنس، فالدنانير أثمان للأشياء في أنفسها.

\subsection{اشتراك العاقلة في دية الخطأ وما دونها}
أجمع أهل العلم وتواترت السنة أن دية قتل الخطأ للمسلم الحر (وهي مائة من الإبل) تتحملها عاقلة الجاني (أي عصبته وقرابته)، وتُقسّط عليهم في ثلاث سنين.
وأجمعوا على أن الجناية العمد تكون في مال الجاني خاصة لا تشاركه عاقلته.

ولكن اختلفوا فيما لو جنى جناية خطأً لا تبلغ دية كاملة، بل تبلغ أقل من الثلث (كالموضحة وهي الشجة التي توضح العظم، وديتها خمس من الإبل). هل تتحملها العاقلة أم تكون في مال الجاني المخطئ وحده؟
قاس الشافعي رحمه الله ذلك قائلاً: إذا كان رسول الله ﷺ قد ألزم العاقلة بتحمل الجناية الكبرى (وهي الدية الكاملة) مواساةً لصاحبهم في الخطأ، فمن باب أولى أو مساواةً أن يتحملوا عنه الجناية الصغرى (كالنصف والثلث وما دونهما) قياساً على الأصل؛ لأنه إذا لزمهم الأكثر لزمهم الأقل في نفس العلة (وهي الخطأ)، وهذا من تمام القياس المطرد.